\section*{The Lumber Camp}\stepcounter{section}\phantomsection\addcontentsline{toc}{section}{The Lumber Camp}
{\entryfont The characters eventually come across a small human lumber camp near the outer reaches of the Woods of Lomond. For generations, the workers here have heard stories of strange lights, elusive forest spirits, and unusual creatures deeper within the woods, but few have ever knowingly encountered the inhabitants of Ty Coeden.

The meeting provides the Caelthir with one of their first opportunities to approach ordinary people of Fife openly and on their own terms. The lumberjacks are wary but not hostile, and how the characters handle this first contact may shape how both groups view one another afterward.}
\subsection*{Approaching the Camp}\stepcounter{subsection}\phantomsection\addcontentsline{toc}{subsection}{Approaching the Camp}
{\entryfont The lumber camp lies in a broad clearing where several recently felled trees have been stripped of their branches and stacked beside rough timber shelters. A handful of wagons stand nearby, while draft horses wait beneath a simple lean-to. The camp is active when the characters arrive, and their approach is noticed almost immediately.

The workers are accustomed to unusual sounds and fleeting shapes within the forest, but openly approaching Fey creatures are another matter entirely. Their first reaction should be one of surprise and apprehension rather than aggression.}
\begin{DndReadAloud}
	The rhythmic crack of axes carries between the trees long before the camp itself comes into view. Ahead, the forest opens into a broad clearing scattered with timber piles, rough shelters, and several heavy wagons loaded with freshly cut logs.

	One by one, the sounds of work begin to stop.

	A lumberjack lowers his axe. Another turns toward the tree line. Within moments, nearly every pair of eyes in the camp is fixed upon you.

	A few workers take cautious steps backward. Others tighten their grip upon tools that suddenly look far more like weapons.

	Finally, a broad-shouldered man near the centre of the camp raises one hand toward the others before calling out:

	"That's close enough. Who - or what - are you?"
\end{DndReadAloud}
\subsection*{First Contact}\stepcounter{subsection}\phantomsection\addcontentsline{toc}{subsection}{First Contact}
{\entryfont The foreman keeps a cautious distance while the other workers watch from behind carts, timber piles, and shelters. Most have spent years hearing stories about strange beings deeper within the Woods of Lomond, and several are visibly unsettled to discover that some of those stories were true.

The characters may introduce themselves however they wish. A calm explanation, an open display of peaceful intent, or a simple gesture of goodwill is enough to prevent the situation from escalating. The lumberjacks are wary, but have no desire to provoke a fight with creatures they barely understand.

Avoid resolving the exchange with a single Charisma Check. Let the characters' words and actions determine how quickly the workers become comfortable around them. Once their peaceful intentions are clear, apprehension gradually gives way to curiosity.\\
\paragraph*{Questions from the Camp} As the workers grow more comfortable, they begin cautiously asking about their unexpected visitors:
\begin{itemize}
	\item "Have you really been living somewhere in these woods all this time?"
	\item "Were those lights people kept seeing yours?"
	\item "Do you know what else lives deeper in there?"
	\item "Why come speak to us now?"
	\item "Are there more of you?"
\end{itemize}
\hfill\\
The characters may reveal as much or as little about Ty Coeden as they consider appropriate. The workers are curious rather than intrusive and generally respect an unwillingness to answer.

If the exchange goes well, the foreman eventually relaxes enough to remark:}
\begin{DndReadAloud}
	"Suppose I always thought the stories were meant to keep children from wandering too far into the trees."

	He glances toward the surrounding forest.

	"Didn't expect the stories to come walking into camp and introduce themselves."
\end{DndReadAloud}
\eject
\subsection*{A Shared Problem}\stepcounter{subsection}\phantomsection\addcontentsline{toc}{subsection}{A Shared Problem}
{\entryfont As the conversation begins to settle into something resembling ordinary hospitality, a sudden crack cuts through the camp.}
\begin{DndReadAloud}
	A sharp crack echoes across the clearing.

	One of the stacked timber piles shifts violently before several heavy logs break loose and roll toward a nearby wagon. Workers shout and scatter as the horses rear against their harnesses.

	Then someone cries out.

	A lumberjack lies pinned beside the wagon, one leg trapped beneath a fallen log while the remaining timber above him creaks dangerously.
	
	The other lumberjacks immediately rush to help, but panic and the unstable timber make coordinated action difficult.
\end{DndReadAloud}
{\noindent\entryfont The incident gives the characters an immediate opportunity to cooperate with the lumberjacks rather than simply speak with them. The problem should be presented as a flexible challenge with several possible solutions. Allow the players to decide how they assist before calling for checks. \\
\paragraph*{Possible Approaches} A successful \textbf{Strength (Athletics) Check} can move or lever the logs, a \textbf{Wisdom (Animal Handling) Check} can calm the horses while suitable spells or tool proficiencies may stabilize the timber. Creative solutions may grant advantage or succeed without a check.\\
\paragraph*{Failure} Failure introduces a complication rather than ending the rescue: another log shifts, the horses begin to bolt, or a character suffers minor damage.\\
\paragraph*{Success} Freeing the worker immediately improves the lumberjacks' disposition toward the characters. Whatever uncertainty remained is tempered by the simple fact that, when something went wrong, both groups acted together. Several workers who previously kept their distance become far more willing to speak openly.}
\subsection*{Around the Camp}\stepcounter{subsection}\phantomsection\addcontentsline{toc}{subsection}{Around the Camp}
{\entryfont With the immediate danger resolved, the atmosphere around the camp becomes noticeably more relaxed. Several workers return to their tasks while others remain nearby to talk, offer food or drink, and satisfy their curiosity about the unexpected visitors.\\
\paragraph*{Possible Topics} The lumberjacks can speak about:
\begin{itemize}
	\item nearby settlements and trade routes
	\item seasonal work and local dangers
	\item life along the edge of the Woods of Lomond
	\item logging practices and the areas they deliberately avoid
\end{itemize}
\hfill\\
The characters may likewise share as much or as little about Ty Coeden as they wish. The scene should remain informal and open-ended, emphasizing that peaceful coexistence can begin through ordinary conversation long before formal agreements or alliances become necessary.}
\subsection*{Parting as Neighbours}\stepcounter{subsection}\phantomsection\addcontentsline{toc}{subsection}{Parting as Neighbours}
{\entryfont Eventually, the characters continue on their way. By now, the lumberjacks' initial fear has softened into cautious familiarity - the Caelthir are no longer merely creatures from old stories, but unfamiliar neighbours. Before they depart, the foreman assures them that peaceful visitors from Ty Coeden will be welcome at the camp. He cannot speak for the wider Kingdom of Fife, but future visits, trade, or cooperation no longer seem impossible.}
\begin{DndReadAloud}
	The foreman gives the group a final, thoughtful look before extending a weathered hand.

	"You've been living beside us longer than any of us knew. Suppose that makes us neighbours, whether we planned on it or not."

	He glances back toward the camp.

	"Come peacefully, and you'll find no trouble here."
\end{DndReadAloud}
{\noindent\entryfont A positive encounter establishes no formal alliance, but gives the characters tangible evidence that peaceful coexistence is possible. When they return to Ty Coeden, they can report that the people beyond the forest are neither guaranteed friends nor the monsters some feared - simply cautious neighbours who may be willing to learn.}